\documentclass[10pt,a4paper]{amsart}
\usepackage[margin=19mm]{geometry}
\usepackage{amsmath,amssymb,mathtools}
\usepackage[scr=rsfs,cal=euler]{mathalfa}
\usepackage{xfrac}
\usepackage[unicode,hidelinks,bookmarks=false]{hyperref}
\newcommand{\ie}{i.e.\ }
\newcommand{\cf}{cf.\ }
\newcommand{\resp}{respectively\ }
\newcommand{\Subsection}{\subsection}
\providecommand{\nobreakdash}{\nobreak\hbox{-}}
\setlength{\emergencystretch}{2em}
\multlinegap0pt
\usepackage{bm}
\usepackage{bbm}
\usepackage{dsfont}
\usepackage{xspace}
\usepackage{cancel}
\usepackage{mathtools}
\usepackage{amsthm}
\makeatletter
\@ifundefined{thm}{\newtheorem{thm}{Thm}}{}
\@ifundefined{lem}{\newtheorem{lem}[thm]{Lemma}}{}
\makeatother
\def\R{{\mathbb R}}
\let\veps=\varepsilon
\title[Stable self-similar singularity formation for infinite energ]{Stable self-similar singularity formation for infinite energy solutions of the incompressible porous medium equations}
\author{Charles Collot, Christophe Prange, Jin Tan}
\date{Source: arXiv:2507.17381v1 (CC BY 4.0)}
\begin{document}
\maketitle

\noindent\emph{Source and licence.} Stable self-similar singularity formation for infinite energy solutions of the incompressible porous medium equations. Version arXiv:2507.17381v1 (CC BY 4.0), distributed under the Creative Commons Attribution 4.0 International licence, as recorded in the arXiv licence metadata for this exact version and shown on the abstract page https://arxiv.org/abs/2507.17381v1. Full rights notice: licenses/arxiv-2507.17381-CC-BY-4.0.txt.\par\smallskip

\noindent\textit{Editorial scope.} Counted unit: the Lem whose source label is \texttt{Le-2.3} in the article (source0.tex lines 741--746, source label \texttt{Le-2.3}), delivered together with the article's own complete original proof of it (source0.tex lines 747--794, one \texttt{proof} environment of 48 lines). No mathematics is written, repaired or re-derived: statement and proof are the article's own text line for line. Required context retained uncounted: eq-L0 (lines 688--690); Le-2.2 (lines 693--706). Every definition, display and label of those lines is retained unchanged. Omitted from this package and not redistributed: 1--687, 691--692, 707--740, 795--1737. Nothing inside a retained excerpt is omitted or altered: every retained body is delivered exactly as the article's lines are written, so no omission receipt of any kind accompanies this package. Citations: the retained excerpts carry no citation command at all, so no bibliography entry of the article is transcribed and no reference is redistributed. Reference adaptation: none was needed and none was made; every reference inside the delivered unit resolves inside it, so no unresolved reference or citation is produced. Printed slips kept literal: no printed word, symbol, hypothesis or number of the counted statement or of its proof was altered. Layout and class: the article's own class is \texttt{amsart} with packages including verbatim, framed, amsfonts, amsthm, amsmath, wasysym, enumerate, fancyhdr, amssymb, chemarrow, mathtools, hyperref, appendix, enumitem; the wrapper is the lane's pinned \texttt{amsart} editorial wrapper at 10pt on A4, which restates the theorem-like environments the retained text needs and adds the article's own macro definitions that the retained text needs (\texttt{\textbackslash R}, \texttt{\textbackslash veps}), with extra wrapper packages (bm, bbm, dsfont, xspace, cancel, mathtools). The delivered numbering is the unit's own. Assets: no retained span contains a figure environment, a graphics-inclusion command, a PDF-page inclusion, a table or a TikZ picture; no image file is copied into this package, the package declares no asset dependency and no image file is redistributed with it. Licence: version arXiv:2507.17381v1 of this article is distributed under the Creative Commons Attribution 4.0 International licence, as recorded in the arXiv licence metadata for this exact version and shown on the abstract page https://arxiv.org/abs/2507.17381v1. A full rights notice is delivered at licenses/arxiv-2507.17381-CC-BY-4.0.txt. Characters: every retained excerpt is pure ASCII, so the delivered engine, XeLaTeX with its default Latin Modern text font, reports no missing character.
\begin{equation}\label{eq-L0}
\partial_t \veps = 2 \cos x \,\veps-\sin x\, \partial_x \veps=L_0 (\veps), \quad \veps(0, x)=\veps_0,\quad {\rm for}~~x\in \Omega^*(t).
 \end{equation}

 \begin{lem}\label{Le-2.2}
 Define $W_{-1}$ as a $2\pi$-periodic function given on $[-\pi, \pi]$ by 
\begin{equation}
{W}_{-1} (x)=\left\{\begin{aligned}
 &|\sin(x) |\sin^2(x/2)\qquad |x|\leq \frac{2\pi}{3},\\
 &\frac{3\sqrt{3}}{8}\qquad\qquad\qquad \quad  x\in [-\pi, - {2\pi}/{3}]\cup [{2\pi}/{3},  \pi].
 \end{aligned}\right.
\end{equation}
Then $W_{-1}\in C^1(\mathbb R)$. For any positive  $W^{1, \infty}(\R)$  function $\widetilde W$ satisfying $\sin x ~{\widetilde W}'(x)\geq 0$ on $\R,$ the 
following holds for $\delta^*>0$ small enough. Let $\veps_0\in  C^{3}(\Omega^*_0)$ such that $\veps_0(x)=O(|x|^{3})$ as $x\to 0$, then $e^{tL0}\veps_0$ the unique global solution of problem \eqref{eq-L0}, satisfies 
\begin{equation}\label{Le2.2-2}
\left\|\frac{e^{tL_0} \veps_0}{W_{-1} \widetilde W}\right\|_{L^\infty{(\Omega^*(t))}}\leq  e^{-t}~
 \left\|\frac{\veps_0}{W_{-1}\widetilde W}\right\|_{L^\infty{(\Omega^*_0)}},\end{equation}
  \end{lem}

 \begin{lem}\label{Le-2.3}
 Let $W_{-1}$ be defined as in Lemma \ref{Le-2.2}.     For any   constant $\theta\in(0, 1)$ and  any  $x^*  \in (-\delta^*, \delta^*),$    if   $\delta^*$ is small enough  then  there exists a positive $W^{1, \infty}(\mathbb R)$ function  $\widetilde{W}_\theta$ with $\sin x \widetilde{W}_\theta'\geq 0$ such that   the following estimate holds:    
\begin{equation}\label{lemeq-3.3}
\left\| \frac{\sin x \int_0^x f(\bar x)\,d\bar x}{W_{-1}\widetilde W_\theta}\right\|_{L^\infty( [-\pi-x^*, \pi-x^*])}\leq ~\theta ~\left\|\frac{f}{W_{-1}\widetilde W_\theta}\right\|_{L^\infty([-\pi-x^*, \pi-x^*])}.
 \end{equation}
 \end{lem}

 \begin{proof}
For any $\theta \in(0, 1)$ and  $C>0$ (to be chosen later),   we define  $\widetilde{W}_\theta$ a 2$\pi-$periodic function whose definition on  $[-\pi, \pi]$ is given  by
 \begin{equation}
 \widetilde{W}_\theta (x) =
 \exp(C|x|/{\theta}),\quad |x|\leq \pi. 
 \end{equation}
    

 Define $v(x)=\frac{f(x)}{W_{-1}\widetilde{W}_{\theta}},$ since $W_{-1} \widetilde{W}_\theta \geq 0$  and both $W_{-1},~\widetilde{W}_\theta$ are even, we only need to discuss about the case when $x\in[0, \pi-x^*]$  and find that 
   \begin{align*}
\left|\frac{\sin x \int_0^x f(\bar x)\,d\bar x}{W_{-1}\widetilde W_\theta}\right|  \leq \|v\|_{L^\infty(\Omega^*)} \frac{ |\sin (x)| \int_0^x W_{-1} \widetilde{W}_\theta \,d\bar x}{W_{-1} \widetilde{W}_\theta}, \quad x\in [0, \pi-x^*].
\end{align*}
 
  
  At first, for $x\in[0,  \frac{2\pi}{3}]$ one has 
 \begin{align*}
\frac{ |\sin (x)| \int_0^x W_{-1} \widetilde{W}_\theta \,d\bar x}{W_{-1} \widetilde{W}_\theta}&\leq  \frac{ \int_0^x \sin^2(\bar x/2)\,\sin \bar x \exp(C\bar x/\theta)  \,d\bar x}{ \sin^2({x/2})\exp(C\theta^{-1} \sin^2(x/2))}  \\
  &\leq  \frac{  \int_0^{x} \exp(C \bar x/\theta)  \,d \,\bar x}{ \exp(Cx\theta)}  \\
&\leq \frac{\theta}{C}.
 \end{align*}
   
Now, we discuss the case  when  $0\leq x^*<\delta^*.$ Thanks to  the fact that $0\leq W_{-1}(x)\leq \frac{3\sqrt{3}}{8}$ for all $x\in\Omega^*$ and $W_{-1}\equiv  \frac{3\sqrt{3}}{8}$ for $x\geq \frac{2\pi}{3},$     we see that  when $x\in[\frac{2\pi}{3}, \pi-x^*],$ 
     \begin{align*}
    \frac{ |\sin (x)| \int_0^x W_{-1} \widetilde{W}_\theta \,d\bar x}{W_{-1} \widetilde{W}_\theta}\leq
\frac{\int_0^x  \exp(C\bar x/\theta)\,d\bar x}{\exp(C x/\theta)}\leq \frac{\theta}{C}.
 \end{align*}
As to the case $-\delta^*<x^*\leq0,$  similarly,  for $x\in[\frac{2\pi}{3}, \pi]$ we can write that
 \begin{align*}
    \frac{ |\sin (x)| \int_0^x W_{-1} \widetilde{W}_\theta \,d\bar x}{W_{-1} \widetilde{W}_\theta}\leq
\frac{\int_0^x  \exp(C\bar x/\theta)\,d\bar x}{\exp(C x/\theta)}\leq \frac{\theta}{C},
 \end{align*}
while for $x\in [\pi, \pi-x^*],$ one has 
  \begin{align*}
      \int_0^x W_{-1} \widetilde{W}_\theta \,d\bar x\leq  \frac{3\sqrt{3}}{8}\left(  \int_0^\pi\exp(C\bar x/\theta)\,d\bar x +  \int_\pi^x \exp(C(2\pi-\bar x)/\theta)\,d\bar x\right),
       \end{align*}
and thus
\begin{align*}
    \frac{ |\sin (x)| \int_0^x W_{-1} \widetilde{W}_\theta \,d\bar x}{W_{-1} \widetilde{W}_\theta}\leq&
 \frac{2\theta \exp(C\pi/\theta)}{C\exp(C(2\pi-x)/\theta)}\\
\leq& \frac{2\theta}{C}\exp(Cx^*/\theta)\leq \frac{6\theta}{C},
 \end{align*}
 provided that $\delta^*\leq \frac{\theta}{C}.$
 Therefore, for all $x\in[0, \pi-x^*]$ and $C$ large enough we have 
 \begin{align*}
\left|\frac{\sin ( x) \int_0^{x}f(\bar x)\,d\bar x}{W_{-1}\widetilde W_\theta}\right|\leq \theta \|v\|_{L^\infty([-\pi-x^*, \pi-x^*])},
 \end{align*}
and there holds \eqref{lemeq-3.3}.
 \end{proof}

\end{document}
